%%%%%%%%%%%%%%%%%%%%%%%%%%%%%%%%%%%%%%%%%%%%%%%%%%%%%
\begin{frame}
    \frametitle{Totales Differential}
    
    \begin{itemize}
        \item Wir definieren das {\bf totale Differential} der Funktion $f(x,y)$ am Punkt $(x_0,y_0)$:
        $$
        df = \frac{\partial f}{\partial x}(x_0,y_0)\,dx +\frac{\partial f}{\partial y}(x_0,y_0)\,dy
        $$
        \item $df$ ist das {\bf totale Differential} am Punkt $(x_0,y_0)$
        \item $dx$ und $dy$ sind infinitesimale Veränderungen der Werte von $x$ und $y$ um den Punkt $(x_0,y_0)$
    \end{itemize}

\end{frame}

%%%%%%%%%%%%%%%%%%%%%%%%%%%%%%%%%%%%%%%%%%%%%%%%%%%%%
\begin{frame}
    \frametitle{Beispiel: Ideales Gasgesetz}

    \begin{itemize}
        \item Ein ideales Gas erfüllt die Bedingung: $p = C\frac{T}{V}$ wo $C$ eine Konstante ist, $p$ der Druck, $V$ das Volumen und $T$ die Temperatur.
        \item Wir möchten Veränderungen des Drucks $p$ bei Veränderungen der anderen Variablen (Temperatur und Volumen) beschreiben um einen Gleichgewichtszustand $(p_0,V_0,T_0)$:
        \begin{equation}
        dp = \frac{\partial p}{\partial T}(T_0,V_0)\,dT +\frac{\partial p}{\partial V}(T_0,V_0)\,dV
        \end{equation}
    \end{itemize}
\end{frame}

%%%%%%%%%%%%%%%%%%%%%%%%%%%%%%%%%%%%%%%%%%%%%%%%%%%%%
\begin{frame}
    \frametitle{Beispiel: Ideales Gasgesetz}

    \begin{itemize}
        \item es gilt:
        $$
        \frac{\partial p}{\partial T}(T_0,V_0) = C\frac{1}{V_0}~~~~~\frac{\partial p}{\partial V}(T_0,V_0) = -C\frac{T_0}{V_0^2}
        $$
        \item daher:
        \begin{equation}
        dp = C\left( \frac{1}{V_0}\cdot dT -\frac{T_0}{V_0^2}\cdot dV\right)
        \end{equation}
        \item Veränderung des Drucks: {\bf positiv} mit Veränderung der Temperatur $T$, {\bf negativ} mit Veränderung des Volumen $V$!
    \end{itemize}
\end{frame}



%%%%%%%%%%%%%%%%%%%%%%%%%%%%%%%%%%%%%%%%%%%%%%%%%%%%%
\begin{frame}
    \frametitle{Linearisierung}

    \begin{itemize}
        \item Anhand der Tangentengleichung oder des totalen Differentials können wir eine {\bf lineare Näherung} für die Funktion $f(x,y)$ am Punkt $(x_0,y_0)$ angeben:
        $$
        f(x,y) \approx f(x_0,y_0) + \frac{\partial f}{\partial x}(x_0,y_0)(x-x_0) + \frac{\partial f}{\partial y}(x_0,y_0)(y-y_0)
        $$
        \item Diese Näherung ist nur gültig, wenn $x$ und $y$ in der Nähe von $x_0$ und $y_0$ liegen!
        \item Beispiel: $f(x,y) = x^2 + y^2$; $x_0=1$, $y_0=1$; $x=1.1$, $y=1.1$
        \item Näherung: $f(1.1,1.1) \approx 2 + 2\cdot(0.1) + 2\cdot(0.1) = 2.4$
        \item reeller Wert: $f(1.1,1.1) = 2.42$
    \end{itemize}
\end{frame}

%%%%%%%%%%%%%%%%%%%%%%%%%%%%%%%%%%%%%%%%%%%%%%%%%%%%%

\begin{frame}
    \frametitle{Linearisierung}

    \begin{itemize}
        \item Wir können die Näherung auch in Form einer Tabelle darstellen:
        \item Beispiel: $f(x,y) = x^2 + y^2$
        \item $L(x,y)=f(x_0,y_0) + \frac{\partial f}{\partial x}(x_0,y_0)(x-x_0) + \frac{\partial f}{\partial y}(x_0,y_0)(y-y_0)=   2 + 2\cdot(x-1) + 2\cdot(y-1)$
    \end{itemize}

\begin{table}[h]
    \begin{tabular}{|c|c|c|c|c|}
    \hline
    $(x,y)$ & Abst. & $f(x,y)$ & $L(x,y)$ & Fehler \\
    \hline
    $(1.0, 1.0)$ & $0.00$ & $2.00$ & $2.00$ & $0.00$ \\
    \hline
    $(1.1, 1.1)$ & $0.14$ & $2.42$ & $2.40$ & $0.02$ \\
    \hline
    $(1.2, 1.2)$ & $0.28$ & $2.88$ & $2.80$ & $0.08$ \\
    \hline
    $(1.3, 1.3)$ & $0.42$ & $3.38$ & $3.20$ & $0.18$ \\
    \hline
    $(1.5, 1.5)$ & $0.71$ & $4.50$ & $4.00$ & $0.50$ \\
    \hline
    $(2.0, 2.0)$ & $1.41$ & $8.00$ & $6.00$ & $2.00$ \\
    \hline
    \end{tabular}
    \end{table}
\end{frame}

%%%%%%%%%%%%%%%%%%%%%%%%%%%%%%%%%%%%%%%%%%%%%%%%%%%%%
\begin{frame}
    \frametitle{Beispiel}


    \begin{itemize}
        \item Um den Gleichgewichtspunkt $(p_0,V_0,T_0)$ kann also die Näherungsformel benutzt werden:
        $$
        \Delta p \approx C\left( \frac{1}{V_0}\cdot \Delta T -\frac{T_0}{V_0^2}\cdot \Delta V\right)
        $$
        \item Beispiel: $V_0=22.4$ l, $T_0=273$ Kelvin
        \item der Druck soll konstant bleiben ($\Delta p=0$), die Temperatur wird um 2 Grad erhöht ($\Delta T=+2$)
        \item Daher ergibt sich eine notwendige Veränderung der Volumens:
        $$
        \Delta V \approx \frac{V_0}{T_0}\cdot \Delta T = +0.16~\text{l}
        $$
 
     \end{itemize}
\end{frame}






%%%%%%%%%%%%%%%%%%%%%%%%%%%%%%%%%%%%%%%%%%%%%%%%%%%%%
\begin{frame}
    \frametitle{Beispiel}

    {\bf Beispiel 2: Näherung einer Funktion}
    \begin{itemize}
        \item Können wir ohne Taschenrechner eine Näherung von $z=\sqrt{2.98^2+4.01^2}$ ausrechnen?
        \item Wir definieren die Funktion $z=f(x,y)= \sqrt{x^2+y^2}$
        \item Wir suchen eine Näherung der Funktion am Punkt $(x_0,y_0,z_0)=(3,4,5)$:
        $$
        z-5 \approx f_x(3,4)(2.98-3) + f_y(3,4)(4.01-4)
        $$
        \item es gilt $f_x = x/\sqrt{x^2+y^2}$ und $f_y = y/\sqrt{x^2+y^2}$
        \item daher $f_x(3,4) = 3/5$ und $f_y(3,4) = 4/5$
        $$
        z\approx 5 + \frac{3}{5}(-0.02) + \frac{4}{5}(0.01) = 5-\frac{2}{500} = 4.996
            $$
    \end{itemize}
\end{frame}

%%%%%%%%%%%%%%%%%%%%%%%%%%%%%%%%%%%%%%%%%%%%%%%%%%%%%
\begin{frame}{Abschlussfragen}

    \begin{enumerate}
        
